% Generated draft: requires mathematical and source review.
\documentclass[11pt]{article}
\usepackage[margin=25mm]{geometry}
\usepackage{fontspec}
\setmainfont{Latin Modern Roman}
\usepackage{amsmath,amssymb,mathtools,unicode-math}
\setmathfont{Latin Modern Math}
\usepackage{enumitem,needspace,xurl,hyperref}
\setlength{\parindent}{0pt}
\setlength{\parskip}{5pt}
\setlength{\emergencystretch}{3em}
\newcommand{\N}{\mathbb{N}}
\newcommand{\Z}{\mathbb{Z}}
\newcommand{\Q}{\mathbb{Q}}
\newcommand{\R}{\mathbb{R}}
\newcommand{\C}{\mathbb{C}}
\newcommand{\F}{\mathbb{F}}
\newcommand{\A}{\mathbb{A}}
\newcommand{\PP}{\mathbb{P}}
\newcommand{\E}{\mathbb{E}}
\newcommand{\eps}{\varepsilon}
\newcommand{\dd}{\,\mathrm{d}}
\begin{document}
\section*{UnsolvedMath ID: 800009}
\textbf{Source problem code:} catalogue entry
\subsection{Definitions and mathematical statement}
Use the standard semigroup signature consisting of one associative binary operation. In this interpretation, a semigroup identity is an equation between nonempty words in variables; a zero element, when present, is not a named constant. The supplied source does not explicitly specify this signature.

A finite semigroup \(S\) is \emph{finitely based} if there is a finite set \(\Sigma\) of identities valid in \(S\) such that every identity valid in \(S\) is an equational consequence of \(\Sigma\). Equivalently, the variety \(\operatorname{var}(S)\) generated by \(S\) has a finite identity basis.

For a finite group \(G\), finite nonempty sets \(I\) and \(\Lambda\), and a matrix \(P=(p_{\lambda i})\) with entries in \(G\cup\{0\}\), suppose that every row and every column of \(P\) contains a nonzero entry. The Rees matrix semigroup
\[
M^0(G;I,\Lambda;P)=\{0\}\cup(I\times G\times\Lambda)
\]
has absorbing element \(0\) and multiplication
\[
(i,g,\lambda)(j,h,\mu)=
\begin{cases}
(i,gp_{\lambda j}h,\mu),&p_{\lambda j}\ne0,\\
0,&p_{\lambda j}=0.
\end{cases}
\]
Under the standard Rees matrix interpretation, these are precisely the finite completely \(0\)-simple semigroups over \(G\), up to isomorphism. Here ``over an Abelian group'' means that this structure group \(G\) is Abelian. Write \(C_2\) for the group of order two.

\textbf{Problem (G. Mashevitzky).} Is there an algorithm which, given finite data \((G,I,\Lambda,P)\) as above with \(G\) Abelian, always terminates and decides whether \(M^0(G;I,\Lambda;P)\) is finitely based? In particular, does such an algorithm exist when the structure group is fixed to be \(C_2\)?

Equivalently, using effective encodings of finite multiplication tables, ask whether the class of finite finitely based completely \(0\)-simple semigroups over Abelian groups is recursive, and ask the same question for the subclass over \(C_2\). This makes explicit the usual interpretation of ``recursive''; the supplied source does not specify an encoding.

The source also reports that an infinite chain of varieties generated by completely \(0\)-simple semigroups over \(C_2\), with finitely based and non-finitely based varieties alternating, was constructed in reference [M4]. This is supplied background, not an independently verified assertion here. The supplied record does not identify [M4] bibliographically or state whether the generating semigroups in that construction are finite.
\subsection{Short English statement}
The question asks whether finite basis existence can be decided for every finite completely \(0\)-simple semigroup whose structure group is Abelian, and specifically when that group has order two. The formulation uses ordinary semigroup identities and finite effective input data; these conventions require review because the source leaves them implicit. The registry labels the problem open, but that status has not been independently verified.
\subsection{Sources}
\begin{itemize}
\item Local UnsolvedMath registry, permanent ID \(800009\), exact source problem code \texttt{catalogue entry}, title \emph{The Lyapin's Notebook --- Problem 9}. The supplied record labels the problem \texttt{open}; this is a supplied status, not an independently verified literature assessment.
\item Bershadsky, Kublanovsky, and Mashevitzky, \emph{The Lyapin's Notebook: Unsolved Problems in Semigroup Theory} (2026), Problem 9, \url{https://arxiv.org/abs/2604.04763}. The supplied record reports access on 2026-10-16 and extraction by \texttt{arxiv\_latex}; neither access nor source contents have been independently checked here. The registry credits these three names in its proposer field and gives proposed year 2026; the extracted problem itself attributes the question to G. Mashevitzky.
\item Reference [M4], cited in the supplied statement for the alternating chain of varieties. No authors, title, publication details, or link for [M4] are supplied, so this reference cannot be identified reliably from the record alone.
\end{itemize}
\end{document}
